\section{Related work}\label{sec:related}

The construction has four properties---finite, solver-attainable, independently checkable, valid against every measurable information-adapted policy---and each of the first three has an established literature. This section places the construction relative to each and states what is and is not claimed.

\paragraph{Viability kernels and two-sided approximation.}
Saint-Pierre's viability kernel algorithm \cite{SaintPierre1994} approximates the finite-horizon kernel by a backward recursion on a grid with a convergence statement for open constraint sets; level-set formulations compute the same objects through Hamilton--Jacobi equations \cite{MitchellBayenTomlin2005}; Lagrangian reformulations in terms of reachable sets scale considerably further \cite{Maidens2013}; and the convex theory of set invariance is surveyed in \cite{Blanchini1999}. \textbf{The two-sided structure is not new, and this paper does not claim it is.} Reachability-based computations produce a guaranteed under-approximation of the kernel together with an over-approximation and use the latter to bound the error of the former \cite[Section~3]{Maidens2013}; the sandwich \eqref{eq:sandwich} is of that kind. What differs is the object and the approximations. Classical schemes discretise a \emph{state space} and approximate a \emph{kernel}; here the discretisation is an outer moment relaxation of information-adapted \emph{controls} together with an inner adversarial selection of labels and scenarios, and the quantity bounded is a safety \emph{value} over an information state. The partial-observation structure, and the requirement that the certificate hold against every measurable policy rather than against one synthesised controller, is what forces the change. Occupation-measure formulations turn optimal control into infinite-dimensional linear programs \cite{Vinter1993} whose LMI relaxations converge \cite{Lasserre2008}, and outer approximations of maximum controlled-invariant sets have been computed this way for polynomial systems \cite{Korda2014}; those are full-information constructions, and their relaxations are semidefinite rather than linear.

\paragraph{Moment relaxations.}
Relaxing a hard problem into convex problems whose values converge monotonically is the moment--SOS programme \cite{Lasserre2001,Parrilo2003}, whose hierarchy produces nondecreasing lower bounds with asymptotic convergence and, importantly, checkable rank and flat-extension conditions under which a relaxation is certified exact. That certification idea is the direct ancestor of the present one, and the differences are concrete: the hierarchy applies to polynomial optimisation over semi-algebraic sets and is solved by semidefinite programming, whereas the relaxations here are \emph{linear} programs over block moments of a linear system; convergence there is in relaxation degree, here along a refinement sequence whose mesh and stored-scenario widths vanish; and an exactness certificate there licenses the extraction of a minimiser, whereas a positive margin here is an \emph{obstruction} certificate.

\paragraph{The scenario approach, and why the guarantee here is not probabilistic.}
Sampling finitely many uncertainty instances and solving the resulting finite program is the scenario approach \cite{CalafioreCampi2005,CalafioreCampi2006}, with sharp sample-complexity results \cite{CampiGaratti2008,CampiGaratti2011,CampiGarattiPrandini2009}. The inner approximation here---finitely many stored labels and scenarios---is superficially the same move but not the same guarantee. The scenario guarantee is probabilistic: violation of the unsampled constraints with probability at most $\epsilon$, with confidence $1-\beta$. The guarantee here is worst-case: scenarios are selected adversarially rather than sampled, the bound holds for every measurable policy, and $\ebar$, $\delta_\beta$ and $Lh_t$ are deterministic inflation certificates, not violation probabilities. (The notational collision on $\beta$ is flagged rather than resolved.)

\paragraph{Minimax reachability, set-membership estimation, and imperfect-information games.}
The quantifier order $\exists$ control $\forall$ disturbance with set-membership uncertainty goes back to \cite{BertsekasRhodes1971}; set-valued state estimation under perturbed observations to \cite{Witsenhausen1968} and the bounding approaches in \cite{Milanese1996}, which are the natural source of the belief cells of Section~\ref{sec:extensions}. In robust linear optimisation, the gap between static and adjustable decisions has a theorem-of-the-alternatives characterisation \cite{WeiZhang2024}; the certificates here measure an information gap of the same nature---a common control forced across indistinguishable branches---but for continuous-time information structures and against measurable policies. Exact algorithms for safety games with imperfect information on finite graphs are available \cite{Chatterjee2007}; what the present paper supplies is the bridge that would make such a finite game a sound abstraction of a continuous system with a measurable policy class. The sufficiency side under imperfect measurement is \cite{Veliov1993}; the necessity side in continuous time is \cite{Abaee2026calc}, whose certificates this paper computes, and whose discrete audits \cite{Abaee2026worked} and measure-dual companion \cite{Abaee2026minimax} are cited where their results are used.

\paragraph{Verified numerics and exact linear programming.}
Enclosing the solution of a continuous problem in rigorously bounded intervals is the programme of interval arithmetic and validated integration; this paper deliberately does not use it. Every exact verification path runs in integer and rational arithmetic and optimality is certified by independent primal and dual witnesses, in the spirit of exact rational LP solution \cite{Applegate2007} but without an exact simplex, because the witnesses are closed-form. The solver is not a proof object. This is a statement about the division of labour, not a claim that interval methods are inferior; they address problems this pipeline does not, and SM~S2 shows where validated enclosures would enter the row budgets.

\paragraph{What is claimed.}
(1) A continuous-to-finite bridge for information states with a two-sided certified bound and convergence along any admissible refinement sequence (Theorem~\ref{thm:bridge}). (2) Dual feasibility as a universal certificate, and an exact atomic dual representation of the value (Corollary~\ref{cor:dual}). (3) Certificate size governed by information--time rank, with a saturating scalar-input construction and a matching depth lower bound (Section~\ref{sec:size}); the negative half---that a single input does not bound the witness count---is the part most likely to surprise. (4) A prescriptive decoding: shared authorities, an exact delay hierarchy, and observer pricing with a certified bracket (Section~\ref{sec:ladder}). (5) Exact verification paths by construction, with the solver demoted to confirmation. Not claimed: the sandwich idea as such, moment relaxation as such, scenario discretisation as such, or interval enclosure.
